\documentclass[11pt]{article}

\usepackage{amsmath,amssymb,amsthm,mathtools}
\usepackage[margin=1in]{geometry}
\usepackage{microtype}
\usepackage[hidelinks]{hyperref}

\newtheorem{theorem}{Theorem}[section]
\newtheorem{lemma}[theorem]{Lemma}
\newtheorem{corollary}[theorem]{Corollary}
\newtheorem{remark}[theorem]{Remark}

\newcommand{\dd}{\,\mathrm{d}}
\newcommand{\one}{\mathbf{1}}
\newcommand{\ip}[2]{\left\langle #1,#2\right\rangle}
\newcommand{\join}{\mathbin{\vee}}

\title{Cones over Forests and the Chromatic-Polynomial Graphon Bound}
\author{}
\date{}

\begin{document}

\maketitle

\begin{abstract}
Let $F$ be a forest with $n$ vertices, $e\geq1$ edges, and $c$ connected
components, so $n=c+e$.  We prove that its cone
$H=K_1\join F$ satisfies
\[
t(H,W)\geq p^c(2p-1)^e=\Phi_H(p)
\]
for every graphon $W$ of edge density $p\geq1/2$.  Since
$\chi(H)=3$, this is the complete required density interval.  The proof
uses the accepted Sidorenko inequality for forests inside the graphon
links and a fully proved weighted rooted-triangle estimate.  It is a
special forest-cone theorem, not a general universal-vertex or join
closure.  In the catalogue on at most six vertices it proves the
previously open Atlas graphs 40, 135, and 136.
\end{abstract}


%%%%%%%%%%%%%%%%%%%%%%%%%%%%%%%%%%%%%%%%%%%%%%%%%%%%%%%%%%%%%%%%%%%%%%%%%%%%%%%
\section{Definitions, chromatic polynomial, and target}
%%%%%%%%%%%%%%%%%%%%%%%%%%%%%%%%%%%%%%%%%%%%%%%%%%%%%%%%%%%%%%%%%%%%%%%%%%%%%%%

Let $(\Omega,\mu)$ be a probability space.  A graphon is a symmetric
measurable function $W\colon\Omega^2\to[0,1]$.  For a finite simple
graph $G$, write
\[
t(G,W)
=
\int_{\Omega^{V(G)}}
\prod_{uv\in E(G)}W(x_u,x_v)
\prod_{u\in V(G)}\dd\mu(x_u).
\]
This is the ordinary, non-injective homomorphism density.  Put
\[
p=t(K_2,W),
\qquad
d(x)=\int_\Omega W(x,y)\dd\mu(y).
\]

Let $F$ be a forest with $n$ vertices, $e\geq1$ edges, and $c$
connected components, including isolated vertices as one-vertex
components.  Each tree component on $v_i$ vertices has chromatic
polynomial $x(x-1)^{v_i-1}$.  Therefore
\[
\chi_F(x)=x^c(x-1)^e,
\qquad
n=c+e.
\]

Define the cone
\[
H=K_1\join F
\]
by adding a new vertex $o$ adjacent to every vertex of $F$.  Colouring
$o$ first gives
\begin{equation}
\chi_H(x)
=x\chi_F(x-1)
=x(x-1)^c(x-2)^e.
\label{eq:chromatic}
\end{equation}
Because $F$ has an edge, $H$ contains a triangle and is $3$-colourable;
hence $\chi(H)=3$.

Let $q=1-p$.  Substitution in the definition of the target yields
\begin{align}
\Phi_H(p)
&=
q^{n+1}
\chi_H\!\left(\frac1q\right)
\notag\\
&=
q^{n+1}\frac1q
\left(\frac pq\right)^c
\left(\frac{2p-1}{q}\right)^e
\notag\\
&=p^c(2p-1)^e,
\label{eq:target}
\end{align}
because $n=c+e$.  The identity holds at $p=1$ by continuous extension.


%%%%%%%%%%%%%%%%%%%%%%%%%%%%%%%%%%%%%%%%%%%%%%%%%%%%%%%%%%%%%%%%%%%%%%%%%%%%%%%
\section{The accepted forest input}
%%%%%%%%%%%%%%%%%%%%%%%%%%%%%%%%%%%%%%%%%%%%%%%%%%%%%%%%%%%%%%%%%%%%%%%%%%%%%%%

We use the following standard, previously accepted positive result.

\begin{theorem}[Sidorenko inequality for forests]
\label{thm:forest-sidorenko}
Let $F$ be a forest with $e$ edges.  Every graphon $V$ of edge density
$z$ satisfies
\[
t(F,V)\geq z^e.
\]
\end{theorem}

Isolated vertices make no change to homomorphism density, disjoint
unions multiply, and the connected case is the classical tree
Sidorenko inequality.  This theorem is the only previously accepted
graph-density result used below.


%%%%%%%%%%%%%%%%%%%%%%%%%%%%%%%%%%%%%%%%%%%%%%%%%%%%%%%%%%%%%%%%%%%%%%%%%%%%%%%
\section{Weighted density in graphon links}
%%%%%%%%%%%%%%%%%%%%%%%%%%%%%%%%%%%%%%%%%%%%%%%%%%%%%%%%%%%%%%%%%%%%%%%%%%%%%%%

Define the rooted triangle density
\[
\tau(x)
=
\int_{\Omega^2}
W(x,y)W(x,z)W(y,z)
\dd\mu(y)\dd\mu(z).
\]

\begin{lemma}[Weighted rooted-triangle inequality]
\label{lem:weighted-triangle}
If $p>0$, then for every integer $h\geq0$,
\[
\int_\Omega d(x)^h\tau(x)\dd\mu(x)
\geq
\frac{2p-1}{p}
\int_\Omega d(x)^{h+2}\dd\mu(x).
\]
\end{lemma}

\begin{proof}
For $j\geq0$, put $M_j=\int_\Omega d(x)^j\dd\mu(x)$.  Let $T_W$
be the integral operator with kernel $W$, and define
\[
I_h=\int_\Omega d(x)^h\tau(x)\dd\mu(x),
\qquad
J_h=\ip{d^h}{T_Wd}.
\]
For $a,b\in[0,1]$, one has $ab\geq a+b-1$.  Apply this to
$W(x,y)$ and $W(x,z)$, multiply by the nonnegative quantity
$d(x)^hW(y,z)$, and integrate.  This gives
\begin{equation}
I_h\geq2J_h-pM_h.
\label{eq:I-J}
\end{equation}

Set $U=1-W$, let $T_U$ be its integral operator, and put
\[
u=T_U\one=1-d,
\qquad q=1-p.
\]
Since $T_Ud=T_U(\one-u)=u-T_Uu$, we have
\[
T_Wd=p\one-T_Ud=d-q\one+T_Uu.
\]
Consequently,
\begin{equation}
J_h=M_{h+1}-qM_h+\ip{d^h}{T_Uu}.
\label{eq:J}
\end{equation}

By symmetry of $U$, twice the difference
\[
\ip{d^h}{T_Uu}-\int_\Omega d(x)^hu(x)^2\dd\mu(x)
\]
equals
\[
\int_{\Omega^2}U(x,y)
\bigl(u(y)-u(x)\bigr)
\bigl(d(x)^h-d(y)^h\bigr)
\dd\mu(x)\dd\mu(y).
\]
This is nonnegative, because $d=1-u$, so $u$ and $d^h$ are oppositely
ordered.  Hence \eqref{eq:J} implies
\begin{align*}
J_h
&\geq
M_{h+1}-qM_h+\int d^h(1-d)^2
\\
&=M_{h+2}-M_{h+1}+pM_h.
\end{align*}
Using this in \eqref{eq:I-J},
\[
I_h\geq2M_{h+2}-2M_{h+1}+pM_h.
\]
Finally,
\begin{align*}
&2M_{h+2}-2M_{h+1}+pM_h
-\frac{2p-1}{p}M_{h+2}
\\
&\qquad=
\frac1p\int_\Omega d(x)^h(d(x)-p)^2\dd\mu(x)
\geq0.
\end{align*}
\end{proof}

For $d(x)>0$, define a probability measure
\[
\dd\mu_x(y)=\frac{W(x,y)}{d(x)}\dd\mu(y)
\]
and let $W_x$ denote the same kernel $W$ on this probability space.
Its edge density is
\[
z_x=t(K_2,W_x)=\frac{\tau(x)}{d(x)^2}.
\]
If $d(x)=0$, set $z_x=0$.

\begin{corollary}[Weighted average link density]
\label{cor:link-average}
For every integer $N\geq2$ and every $p>0$,
\[
\frac{\int_\Omega d(x)^Nz_x\dd\mu(x)}
{\int_\Omega d(x)^N\dd\mu(x)}
\geq2-\frac1p.
\]
\end{corollary}

\begin{proof}
The numerator is $\int d^{N-2}\tau$.  Apply
Lemma~\ref{lem:weighted-triangle} with $h=N-2$.  The denominator is
positive because $p>0$.
\end{proof}


%%%%%%%%%%%%%%%%%%%%%%%%%%%%%%%%%%%%%%%%%%%%%%%%%%%%%%%%%%%%%%%%%%%%%%%%%%%%%%%
\section{The forest-cone theorem}
%%%%%%%%%%%%%%%%%%%%%%%%%%%%%%%%%%%%%%%%%%%%%%%%%%%%%%%%%%%%%%%%%%%%%%%%%%%%%%%

\begin{theorem}[Cones over forests]
\label{thm:main}
Let $F$ be a forest with $n$ vertices, $e\geq1$ edges, and $c$
connected components.  Let $H=K_1\join F$.  Every graphon $W$ with
edge density $p\geq1/2$ satisfies
\[
\boxed{
t(H,W)
\geq p^c(2p-1)^e
=\Phi_H(p).
}
\]
Since $\chi(H)=3$, the interval $p\geq1/2$ is exactly the full required
interval $p\geq1-1/(\chi(H)-1)$.
\end{theorem}

\begin{proof}
Condition on the image $x$ of the cone vertex $o$.  Directly from the
homomorphism integral,
\begin{equation}
t(H,W)
=\int_\Omega d(x)^n t(F,W_x)\dd\mu(x).
\label{eq:link-identity}
\end{equation}
At points where $d(x)=0$, both the original conditional integral and
the displayed integrand vanish.  Thus no pointwise positive-degree
assumption is being made.

By Theorem~\ref{thm:forest-sidorenko},
\[
t(F,W_x)\geq z_x^e.
\]
Put $M_n=\int d^n$.  Since $e\geq1$, the function $z\mapsto z^e$ is
convex and nondecreasing on $[0,1]$.  Jensen's inequality with respect
to the probability measure $d(x)^n\dd\mu(x)/M_n$, followed by
Corollary~\ref{cor:link-average}, gives
\begin{align*}
t(H,W)
&\geq\int_\Omega d(x)^nz_x^e\dd\mu(x)
\\
&\geq
M_n
\left(
\frac{\int d(x)^nz_x\dd\mu(x)}{M_n}
\right)^e
\\
&\geq
M_n\left(2-\frac1p\right)^e.
\end{align*}
The last step uses $p\geq1/2$, so the lower bound on the weighted mean
is nonnegative.  A second application of Jensen gives $M_n\geq p^n$.
Therefore
\begin{align*}
t(H,W)
&\geq
p^n\left(\frac{2p-1}{p}\right)^e
\\
&=p^{n-e}(2p-1)^e
=p^c(2p-1)^e.
\end{align*}
This is exactly \eqref{eq:target}.
\end{proof}


%%%%%%%%%%%%%%%%%%%%%%%%%%%%%%%%%%%%%%%%%%%%%%%%%%%%%%%%%%%%%%%%%%%%%%%%%%%%%%%
\section{Endpoints and equality examples}
%%%%%%%%%%%%%%%%%%%%%%%%%%%%%%%%%%%%%%%%%%%%%%%%%%%%%%%%%%%%%%%%%%%%%%%%%%%%%%%

At $p=1/2$, the target is zero because $e\geq1$, so the conclusion
also follows directly from nonnegativity.  The balanced complete
bipartite graphon has no triangle and attains equality.  At $p=1$, one
has $W=1$ almost everywhere, and both sides equal $1$.

For every integer $k\geq2$, the balanced complete $k$-partite graphon
$T_k$ has edge density $p=1-1/k$ and satisfies
\[
t(H,T_k)=\frac{\chi_H(k)}{k^{v(H)}}=\Phi_H(p).
\]
Thus the bound has the expected equality examples at every multipartite
grid point in the required interval.


%%%%%%%%%%%%%%%%%%%%%%%%%%%%%%%%%%%%%%%%%%%%%%%%%%%%%%%%%%%%%%%%%%%%%%%%%%%%%%%
\section{Catalogue consequences}
%%%%%%%%%%%%%%%%%%%%%%%%%%%%%%%%%%%%%%%%%%%%%%%%%%%%%%%%%%%%%%%%%%%%%%%%%%%%%%%

\begin{corollary}[New positive Atlas cases]
The following previously open graphs satisfy the desired inequality for
every graphon throughout the full interval $p\geq1/2$:

\begin{center}
\begin{tabular}{cclll}
\hline
Atlas ID & graph6 & $(c,e)$ for $F$ & $\chi_H(x)$ & $\Phi_H(p)$ \\
\hline
40  & \texttt{DjW}  & $(2,2)$ & $x(x-1)^2(x-2)^2$ & $p^2(2p-1)^2$ \\
135 & \texttt{EzW\_} & $(2,3)$ & $x(x-1)^2(x-2)^3$ & $p^2(2p-1)^3$ \\
136 & \texttt{Ejt?} & $(2,3)$ & $x(x-1)^2(x-2)^3$ & $p^2(2p-1)^3$ \\
\hline
\end{tabular}
\end{center}
\end{corollary}

\begin{proof}
Deleting the universal vertex from Atlas 40 leaves a three-vertex path
and one isolated vertex, so $(c,e)=(2,2)$.  Deleting the universal
vertex from either Atlas 135 or Atlas 136 leaves a four-vertex tree and
one isolated vertex, so $(c,e)=(2,3)$.  The two four-vertex trees are
nonisomorphic, but Theorem~\ref{thm:main} applies to every forest and
depends only on $c$ and $e$.  The graph6 codes are independent
machine-readable identifiers.
\end{proof}


%%%%%%%%%%%%%%%%%%%%%%%%%%%%%%%%%%%%%%%%%%%%%%%%%%%%%%%%%%%%%%%%%%%%%%%%%%%%%%%
\section{Scope and accepted input}
%%%%%%%%%%%%%%%%%%%%%%%%%%%%%%%%%%%%%%%%%%%%%%%%%%%%%%%%%%%%%%%%%%%%%%%%%%%%%%%

The only previously accepted graph-density result used is the forest
Sidorenko inequality, Theorem~\ref{thm:forest-sidorenko}.  The weighted
link estimate is proved here in full.

The theorem does not say that $K_1\join G$ is positive whenever $G$ is
positive.  Its proof depends specifically on the forest inequality
$t(F,V)\geq z^{e(F)}$ and on the convexity of the monomial $z^e$.
Consequently, it establishes exactly cones over forests and no general
universal-vertex or arbitrary-join closure.

The argument applies directly to arbitrary graphons and ordinary
homomorphism densities.  It uses no finite-step approximation,
regularity, injectivity, minimum-degree assumption, or pointwise lower
bound on the degree function.

%%%%%%%%%%%%%%%%%%%%%%%%%%%%%%%%%%%%%%%%%%%%%%%%%%%%%%%%%%%%%%%%%%%%%%%%%%%%%%%
\appendix
\section{What the Lean formalization actually proves}
%%%%%%%%%%%%%%%%%%%%%%%%%%%%%%%%%%%%%%%%%%%%%%%%%%%%%%%%%%%%%%%%%%%%%%%%%%%%%%%

The development is \texttt{lean/Taeyoung/Methods/ForestCone/Rows.lean}, over
\texttt{Methods/ConeBound.lean}.  All five Atlas rows this methodology is
credited with --- $40$, $111$, $117$, $135$, $136$ --- are
\texttt{verified}, as is $192$, the cone over $40$.  One difference of
substance.

\subsection*{Theorem~\ref{thm:forest-sidorenko} is never assumed}

The note takes the forest Sidorenko inequality as an accepted input.  Nothing
in the catalogue may rest on an accepted input, so Theorem~\ref{thm:main} is
formalized in the conditional form
\texttt{satisfiesLowerBound\_coneGraph\_of\_sidorenko}, with $t(F,V)\geq z^e$
a hypothesis, and that hypothesis is discharged separately for each row.  Four
of the five rows turn out not to need anything of the kind: deleting the
universal vertex leaves

\begin{center}
\begin{tabular}{r|l|l}
Atlas & base $F$ & why $t(F,V)\geq z^{e(F)}$\\\hline
$40$ & $K_{1,2}\sqcup K_1$ & $t(K_{1,n},V)=M_n\geq z^n$, Jensen\\
$111$ & $K_{1,2}\sqcup2K_1$ & the same\\
$117$ & $K_2\sqcup K_2\sqcup K_1$ & multiplicativity, with equality\\
$135$ & $K_{1,3}\sqcup K_1$ & Jensen again\\
$136$ & $P_4\sqcup K_1$ & \textbf{not} a star: needs $t(P_4,V)\geq z^3$
\end{tabular}
\end{center}

so only Atlas $136$ uses a genuine tree inequality, and only the single case
$P_4$.  That case is proved outright in \texttt{Methods/PathSidorenko.lean}, by
three-factor H\"older on the almost-everywhere factorization
$W=(Wdd')^{1/3}(W/d)^{1/3}(W/d')^{1/3}$.  Tree Sidorenko in general, and hence
Theorem~\ref{thm:forest-sidorenko} in general, is not formalized and is not
needed.

Lemma~\ref{lem:weighted-triangle} is formalized as stated, and is shared with
the clique common-leaf and cone families.

\end{document}
